% Part: normal-modal-logic
% Chapter: filtrations
% Section: euclidean-filtrations

\documentclass[../../../include/open-logic-section]{subfiles}

\begin{document}

\olfileid[fr]{nml}{fil}{euc}

\olsection{Filtrations de modèles euclidiens}

La méthode de la \olref[acc]{sec} ne fonctionne pas pour les modèles
euclidiens, ni pour les modèles à la fois sériels et euclidiens.
Considérons le modèle en haut de la \olref{fig:ser-eucl} : il est à la
fois euclidien et sériel. Soit $\Gamma = \{p, \Box p \}$. Dans une
filtration par $\Gamma$, on a $[w_1] = [w_3]$, car $w_1$ et $w_3$
sont les seuls mondes qui s'accordent sur~$\Gamma$. Toute filtration
doit aussi conserver les flèches héritées de $\mModel{M}$, comme dans la
\olref{fig:ser-eucl2}. Les classes $[w_2]$ et $[w_4]$ ont le prédécesseur
commun $[w_1]$. Dans une relation euclidienne, chacune
devrait donc accéder à l'autre. En particulier, $[w_4]$
accéderait à $[w_2]$ et à $[w_5]$, de sorte que $[w_2]$
devrait accéder à $[w_5]$. Or la filtration doit rendre
$\Box p$ vraie en~$[w_2]$ et $p$ fausse en~$[w_5]$.
Cette flèche est donc interdite par la condition de filtration.
Aucune relation euclidienne ne convient.

\begin{figure}[htpb]
  \centering
  \begin{tikzpicture}[modal]
    \node[world] (w1)
          [label={left: \mFalse{p}},
            label={below: $\mSat{{}}{\Box p}$}] {$w_1$};
    \node[world] (w2)
          [label={right: \mTrue{p}},
            label = {below: $\mSat{{}}{\Box p}$},
            right=of w1] {$w_2$};
    \draw[->] (w1) to (w2);
    \draw[reflexive above] (w2) to (w2);
    \node[world] (w3)
          [label={left: \mFalse{p}},
            label={below: $\mSat{{}}{\Box p}$},
            below=of w1] {$w_3$};
    \node[world] (w4)
          [label={right:\mTrue{p}},
            label={below: $\mSat/{{}}{\Box p}$},
            right=of w3] {$w_4$};
    \draw[->] (w3) to (w4);
    \draw[reflexive above] (w4) to (w4);
    \node[world] (w5)
          [label={right: \mFalse{p}},
            label=below:$\mSat/{{}}{\Box p}$,
            right=of w4] {$w_5$};
    \draw[->, bend left] (w4) to (w5);
    \draw[->, bend left] (w5) to (w4);
    \draw[reflexive above] (w5) to (w5);
  \end{tikzpicture}
  \caption{Un modèle sériel et euclidien.}
  \ollabel{fig:ser-eucl}
\end{figure}

\begin{figure}[ht]
  \centering
  \begin{tikzpicture}[modal,every text node part/.style={align=center}]
    \node[world] (w1)
          [label={left: \mFalse{p}},
            label={above left:$[w_1]=[w_3]$},
            label={below: $\mSat{{}}{\Box p}$}] {$[w_1]$};
    \node[world] (w2)
         [label={right: \mTrue{p}},
           label = {below: $\mSat{{}}{\Box p}$},
           right=of w1, yshift=1.5cm] {$[w_2]$};
    \draw[->] (w1) to (w2);
    \draw[reflexive above] (w2) to (w2);
    \node[world] (w4)
         [label={right:\mTrue{p}},
           label={below: $\mSat/{{}}{\Box p}$},
           right=of w1,yshift=-1.5cm] {$[w_4]$};
    \draw[->] (w1) to (w4);
    \draw[reflexive above] (w4) to (w4);
    \node[world] (w5)
         [label={right: \mFalse{p}},
             label=below:$\mSat/{{}}{\Box p}$,
             right=of w4] {$[w_5]$};
    \draw[->, bend left] (w4) to (w5);
    \draw[->, bend left] (w5) to (w4);
    \draw[reflexive above] (w5) to (w5);
  \end{tikzpicture}
  \caption{La filtration du modèle de la \olref{fig:ser-eucl}.}
  \ollabel{fig:ser-eucl2}
\end{figure}

En particulier, pour obtenir une filtration euclidienne, il ne suffit
pas de filtrer par un ensemble $\Gamma$ quelconque clos par
sous-!!{formula}s. Il faut plutôt considérer des ensembles $\Gamma$
\emph{modalement clos} (voir la \olref[pre]{defn:modallyclosed}). Ces
ensembles de formules, s’ils sont non vides, sont infinis et ne donnent donc pas immédiatement
la propriété du modèle fini ni la décidabilité du système correspondant.

\begin{thm}\ollabel{thm:modal-closed-filt}
  Supposons $\Gamma$ modalement clos,
  $\mModel{M}=\tuple{W,R,V}$, et
  $\mModel{M^*} = \tuple{W^*,R^*,V^*}$ une filtration la plus grossière
  de $\mModel{M}$.
  \begin{enumerate}
  \item Si $\mModel{M}$ est symétrique, $\mModel{M^*}$ l'est aussi.
  \item Si $\mModel{M}$ est transitif, $\mModel{M^*}$ l'est aussi.
  \item Si $\mModel{M}$ est euclidien, $\mModel{M^*}$ l'est aussi.
  \end{enumerate}
\end{thm}

\begin{proof}
\begin{enumerate}
  \item Exercice. Utilisez le fait que \Ax{B} et $\Ax{B_\Diamond}$
    sont valides dans tous les modèles symétriques.
  \item Puisque $\mModel{M^*}$ est une filtration la plus grossière,
    on a par définition $R^*[u][v]$ si et seulement si $C_1(u,v)$.
    Pour la transitivité, supposons $C_1(u,v)$ et $C_1(v,w)$ ; il
    faut montrer $C_1(u,w)$. \iftag{prvBox}{Supposons
      $\mSat{M}{\Box !A}[u]$ ; alors
      $\mSat{M}{\Box\Box!A}[u]$ puisque \Ax{4} est valide dans tous
      les modèles transitifs. Comme $\Box\Box!A \in \Gamma$ par
      clôture modale, $C_1(u,v)$ donne
      $\mSat{M}{\Box!A}[v]$, puis $C_1(v,w)$ donne
      $\mSat{M}{!A}[w]$. }{}\iftag{prvDiamond}{Supposons
      $\mSat{M}{!A}[w]$ ; alors
      $\mSat{M}{\Diamond !A}[v]$ par $C_1(v,w)$, puisque
      $\Diamond !A \in \Gamma$ par clôture modale. Par $C_1(u,v)$,
      on obtient $\mSat{M}{\Diamond\Diamond !A}[u]$, puisque
      $\Diamond\Diamond!A \in \Gamma$ par clôture modale. Comme
      $\Ax{4}_\Diamond$ est valide dans tous les modèles transitifs,
      $\mSat{M}{\Diamond!A}[u]$.}{}
  \item Exercice. Utilisez le fait que \Ax{5} et $\Ax{5_\Diamond}$
    sont valides dans tous les modèles euclidiens.
\end{enumerate}
\end{proof}

\begin{prob}
  Complétez la démonstration du \olref[nml][fil][euc]{thm:modal-closed-filt}.
\end{prob}
\begin{explain}
Dans la \olref{fig:ser-eucl}, le monde~$w_2$ a pour seul
successeur lui-même. Ce successeur suffit à satisfaire la
sérialité en~$w_2$ et impose sa boucle par la propriété
euclidienne appliquée à la flèche de~$w_1$ vers~$w_2$.
La flèche héritée figure aussi sur la classe de~$w_2$
dans la \olref{fig:ser-eucl2}.
\end{explain}
\end{document}
